\documentclass[a4paper,11pt]{amsart}

% =========================================================================
%  Informe de laboratorio — Práctica 2: Ondas estacionarias
%  Estilo: el del informe del efecto Doppler (artículo científico
%  minimalista, Palatino, versalitas, enunciados numerados por sección,
%  paleta de azules y grises), con los encabezados de sección, el índice y
%  la paginación inspirados en el guion de Cosmosofía.
% =========================================================================

% --- IDIOMA Y CODIFICACIÓN ------------------------------------------------
\usepackage[utf8]{inputenc}
\usepackage[T1]{fontenc}
\usepackage[spanish,es-tabla,es-nodecimaldot]{babel}

% --- TIPOGRAFÍA (Palatino) --------------------------------------------------
\usepackage{mathpazo}
\linespread{1.05}
\usepackage[final]{microtype}

% --- MATEMÁTICAS Y UNIDADES ------------------------------------------------
\usepackage{amsmath, amssymb, mathtools}
\usepackage{bm}
\usepackage{siunitx}
\sisetup{
  output-decimal-marker = {,},
  group-minimum-digits  = 4,
  group-digits          = integer,
  group-separator       = {\,},
  per-mode              = symbol,
  separate-uncertainty  = true,
  range-phrase          = {\text{ a }},
  detect-all            = false
}

% --- PÁGINA, COLORES Y GRÁFICOS --------------------------------------------
\usepackage[a4paper,margin=2.6cm,footskip=1.45cm]{geometry}
\usepackage{graphicx}
\usepackage{xcolor}
\definecolor{azulprof}{HTML}{1B365D}    % azul marino
\definecolor{azulacero}{HTML}{3B6EA8}   % azul acero
\definecolor{azulclaro}{HTML}{B7CBE3}   % azul hielo
\definecolor{grispizarra}{HTML}{5B6770} % gris pizarra
\definecolor{grisclaro}{HTML}{E4E8EC}   % gris perla
\definecolor{azulfondo}{HTML}{F3F7FB}   % fondo tenue de referencias y resultados
\definecolor{acento}{HTML}{B8800A}      % acento sobrio (ámbar oscuro) para comparaciones
\usepackage{tikz}
\usetikzlibrary{babel, arrows.meta, calc, positioning, decorations.pathmorphing,
                decorations.markings, patterns}
\usepackage{pgfplots}
\usepackage{pgfplotstable}
\usepgfplotslibrary{groupplots}
\pgfplotsset{compat=1.17}
\usepackage{booktabs}
\usepackage{multirow}
\usepackage[most]{tcolorbox}
\usepackage{array}
\usepackage{enumitem}
\usepackage{needspace}
\setlist[itemize]{leftmargin=1.6em, itemsep=2pt, topsep=4pt}
\setlist[enumerate]{leftmargin=1.8em, itemsep=2pt, topsep=4pt}

% --- ENLACES Y REFERENCIAS CRUZADAS ----------------------------------------
\usepackage[colorlinks=true, linkcolor=azulacero, citecolor=azulacero,
            urlcolor=azulacero, plainpages=false, pdfusetitle=false,
            pdftitle={Práctica 2: Ondas estacionarias},
            pdfauthor={David Gómez García; Daniel Martín Guerrero}]{hyperref}
\usepackage[spanish,nameinlink]{cleveref}

% --- BIBLIOGRAFÍA -----------------------------------------------------------
\usepackage{csquotes}
\usepackage[backend=biber, style=numeric, sorting=none, giveninits=true,
            maxbibnames=6, minbibnames=3, doi=true, url=true,
            isbn=false]{biblatex}
\addbibresource{referencias.bib}
\DeclareFieldFormat[article,techreport,unpublished]{title}{#1\isdot}
\DeclareFieldFormat{url}{\url{#1}}
\renewcommand*{\bibfont}{\small}
\setlength{\bibitemsep}{3pt}


% =========================================================================
%        CABECERAS, PIE DE PÁGINA, SECCIONES E ÍNDICE (aire de Cosmosofía)
% =========================================================================
% Pequeña espiral: viñeta de la paginación y ornamento entre secciones
\newcommand{\espiralita}[1][0.9]{%
  \tikz[baseline=-0.6ex, scale=#1]{%
    \draw[azulacero, line width=0.45pt]
      plot[domain=0:15.7, samples=70, variable=\t, smooth]
      ({0.0105*\t*cos(\t r)},{0.0105*\t*sin(\t r)});
    \fill[azulprof] (0,0) circle (0.55pt);}}
\newcommand{\ornamento}{%
  \par\vspace{10pt}{\centering
  \tikz{\draw[azulclaro, line width=0.4pt] (-1.6,0) -- (-0.32,0);
        \draw[azulclaro, line width=0.4pt] (0.32,0) -- (1.6,0);
        \node at (0,0) {\espiralita[1.6]};}\par}\vspace{6pt}}

\makeatletter
\def\ps@headings{%
  \let\@mkboth\markboth
  \def\@evenhead{\normalfont\footnotesize\color{grispizarra}\scshape
    \textls[60]{Laboratorio de Mecánica y Ondas}\hfil
    \textls[60]{Ondas estacionarias en una cuerda}}%
  \let\@oddhead\@evenhead
  \def\@evenfoot{\hfil\normalfont\small\color{azulacero}%
    \espiralita\enspace\thepage\enspace\espiralita\hfil}%
  \let\@oddfoot\@evenfoot}
\def\ps@firstpage{\ps@plain\def\@oddhead{}\let\@evenhead\@oddhead
  \def\@oddfoot{\hfil\normalfont\small\color{azulacero}%
    \espiralita\enspace\thepage\enspace\espiralita\hfil}%
  \let\@evenfoot\@oddfoot}
\pagestyle{headings}

% Secciones: número en azul acero, título en versalitas espaciadas; debajo,
% un subtítulo en cursiva y un filete rematado en un pequeño círculo.
\def\section{\@startsection{section}{1}%
  \z@{1.3\linespacing\@plus.8\linespacing}{.25\linespacing}%
  {\normalfont\Large\color{azulprof}\scshape\lsstyle\raggedright}}
\let\@seccntformatorig\@seccntformat
\def\@seccntformat#1{%
  \ifnum\pdfstrcmp{#1}{section}=\z@
    {\color{azulacero}\upshape\csname the#1\endcsname}\hspace{0.55em}%
  \else\@seccntformatorig{#1}\fi}
\newcommand{\filete}{%
  \tikz[baseline]{\draw[azulclaro, line width=0.45pt] (0,0) -- (\linewidth-5pt,0);
    \draw[azulacero, line width=0.5pt, fill=white] (\linewidth-2.6pt,0) circle (2.3pt);}}
\DeclareRobustCommand{\tocsub}[1]{%
  \ \textperiodcentered\ {\itshape\color{grispizarra}#1}}
\newcommand{\subtitulo}[1]{%
  \par\nopagebreak\vspace{-2pt}%
  \noindent{\itshape\color{grispizarra}#1}\par\nopagebreak\vspace{1pt}%
  \noindent\filete\par\nopagebreak\vspace{0.55\baselineskip}}
% \seccion{Título}{Subtítulo}: sección numerada con subtítulo (índice incluido).
% amsart escribe en el índice el título principal, no el opcional, así que el
% subtítulo se añade a esa entrada mediante \@tocextra.
\let\@tocextra\@empty
\def\@tocwriteb#1#2#3{%
  \begingroup
    \def\@tocline##1##2##3##4##5##6{%
      \ifnum##1>\c@tocdepth
      \else \sbox\z@{##5\let\indentlabel\@tochangmeasure##6}\fi}%
    \csname l@#2\endcsname{#1{\csname#2name\endcsname}{\@secnumber}{}}%
  \endgroup
  \addcontentsline{toc}{#2}%
    {\protect#1{\csname#2name\endcsname}{\@secnumber}{#3\@tocextra}}%
  \global\let\@tocextra\@empty}
\newcommand{\seccion}[2]{%
  \gdef\@tocextra{\tocsub{#2}}%
  \section{#1}\subtitulo{#2}}

% Índice: número en azul acero, título en azul marino y subtítulo en cursiva
\def\l@section{\@tocline{1}{7pt}{0pt}{}{}}
\renewcommand{\tocsection}[3]{%
  \indentlabel{\@ifnotempty{#2}{{\color{azulacero}#2}\quad}}{\color{azulprof}#3}}
\def\l@subsection{\@tocline{2}{1pt}{1.55pc}{2.4pc}{\small}}
\renewcommand{\contentsnamefont}{\scshape\color{azulprof}\lsstyle}
\makeatother
\numberwithin{equation}{section}
\setcounter{tocdepth}{2}
\makeatletter
\def\@@and{y}
\renewcommand{\abstractname}{\color{azulacero}\lsstyle Resumen}
\makeatother

% =========================================================================
%               NOTACIÓN Y MACROS DE USO FRECUENTE
% =========================================================================
\newcommand{\mub}{\mu_{\mathrm{b}}}          % densidad lineal medida con la báscula
\newcommand{\mus}{\mu_{\mathrm{s}}}          % densidad deducida de una pendiente
\newcommand{\pd}[2]{\frac{\partial #1}{\partial #2}}
\newcommand{\pdd}[2]{\frac{\partial^{2} #1}{\partial #2^{2}}}
\newcommand{\dif}{\mathrm{d}}
\newcommand{\ddt}[1]{\frac{\dif #1}{\dif t}}
\newcommand{\vect}[1]{\vec{#1}}
\newcommand{\Tcoef}{t_{95,\nu}}               % factor de cobertura al 95 %
\newcommand{\ext}[1]{U(#1)}                   % incertidumbre ampliada
\newcommand{\mean}[1]{\overline{#1}}
% Dato pendiente de incorporar (se resalta con discreción)
\newcommand{\pend}[1]{\textcolor{grispizarra}{\colorbox{grisclaro}{\strut\,#1\,}}}
\newcommand{\pendiente}[1]{\textcolor{grispizarra}{\textit{[#1]}}}
% Hueco reservado para una gráfica o figura que aún no existe
\newcommand{\huecofigura}[2][5.2cm]{%
  \tikz{\draw[grispizarra!45, dashed, line width=0.5pt, rounded corners=2pt, fill=azulfondo]
          (0,0) rectangle (0.92\linewidth,#1);
        \node[text=grispizarra, font=\small\itshape, align=center, text width=0.8\linewidth]
          at (0.46\linewidth,#1/2) {#2};}}

% =========================================================================
%        ENTORNOS FORMALES (enunciados numerados, estilo artículo)
% =========================================================================
\usepackage{amsthm}

% Enunciados en cursiva (resultados generales)
\newtheoremstyle{cienplano}{10pt}{10pt}{\itshape}{}%
  {\bfseries\color{azulprof}}{.}{0.6em}%
  {\thmname{#1}\thmnumber{ #2}\thmnote{\normalfont\upshape\ (#3)}}
% Definiciones, cálculos y estimaciones (cuerpo en redonda)
\newtheoremstyle{ciendef}{10pt}{10pt}{\normalfont}{}%
  {\bfseries\color{azulprof}}{.}{0.6em}%
  {\thmname{#1}\thmnumber{ #2}\thmnote{\normalfont\upshape\ (#3)}}
% Observaciones (encabezado en cursiva)
\newtheoremstyle{cienobs}{8pt}{8pt}{\normalfont}{}%
  {\itshape\color{azulprof}}{.}{0.6em}%
  {\thmname{#1}\thmnumber{ #2}\thmnote{\normalfont\upshape\ (#3)}}

\theoremstyle{cienplano}
\newtheorem{proposicion}{Proposición}[section]
\newtheorem{teorema}[proposicion]{Teorema}
\newtheorem{resultado}[proposicion]{Resultado}
\newtheorem{conclusion}[proposicion]{Conclusión}
\theoremstyle{ciendef}
\newtheorem{definicion}[proposicion]{Definición}
\newtheorem{hipotesis}[proposicion]{Hipótesis}
\newtheorem{calculo}[proposicion]{Cálculo de incertidumbre}
\newtheorem{estimacion}[proposicion]{Estimación lineal}
\theoremstyle{cienobs}
\newtheorem{observacion}[proposicion]{Observación}
\newtheorem*{demostracion}{Demostración}

% --- Referencia formalizada --------------------------------------------------
% Caja fina con fondo tenue y título sobre el borde. Cuerpo en cursiva y, al
% pie, la fuente completa con autor, título y número de cita. Comparte
% contador con los enunciados.
% Uso: \begin{referencia}[título]{clave-bib} ... \end{referencia}
\newcounter{referencia}
\renewcommand{\thereferencia}{\theproposicion}
\newtcolorbox{referenciabox}[2]{%
  enhanced, breakable, sharp corners,
  colback=azulfondo, colframe=azulclaro, boxrule=0.4pt,
  left=10pt, right=10pt, top=7pt, bottom=6pt,
  before skip=14pt, after skip=12pt,
  attach boxed title to top left={yshift*=-\tcboxedtitleheight/2, xshift=8pt},
  boxed title style={colback=white, colframe=white, sharp corners, boxrule=0pt,
                     left=3pt, right=3pt, top=1pt, bottom=1pt},
  fonttitle=\bfseries\color{azulprof}, coltitle=azulprof,
  fontupper=\itshape,
  title={Referencia~\thereferencia\ifx\relax#1\relax\else\ \normalfont(#1)\fi},
  after upper={\par\vspace{5pt}{\color{azulclaro}\hrule height 0.4pt}\vspace{4pt}%
    {\upshape\footnotesize\color{grispizarra}\raggedleft
     \citeauthor{#2}, \citetitle{#2}~\cite{#2}.\par}}}
\newenvironment{referencia}[2][]
  {\par\stepcounter{proposicion}\refstepcounter{referencia}\begin{referenciabox}{#1}{#2}}
  {\end{referenciabox}}
\crefname{referencia}{referencia}{referencias}

% --- Método (caja azulada para los fundamentos del tratamiento de datos) -----
\newcounter{metodo}
\renewcommand{\themetodo}{\theproposicion}
\newtcolorbox{metodobox}[1]{%
  enhanced, breakable, sharp corners,
  colback=azulfondo, colframe=azulclaro, boxrule=0.4pt,
  left=10pt, right=10pt, top=7pt, bottom=5pt,
  before skip=14pt, after skip=12pt,
  attach boxed title to top left={yshift*=-\tcboxedtitleheight/2, xshift=8pt},
  boxed title style={colback=white, colframe=white, sharp corners, boxrule=0pt,
                     left=3pt, right=3pt, top=1pt, bottom=1pt},
  fonttitle=\bfseries\color{azulprof}, coltitle=azulprof,
  title={Método~\themetodo\ifx\relax#1\relax\else\ \normalfont(#1)\fi}}
\newenvironment{metodo}[1][]
  {\par\stepcounter{proposicion}\refstepcounter{metodo}\begin{metodobox}{#1}}
  {\end{metodobox}}
\crefname{metodo}{método}{métodos}
\Crefname{metodo}{Método}{Métodos}
\Crefname{referencia}{Referencia}{Referencias}

% --- Resultado destacado ----------------------------------------------------
% Fondo azulado tenue y filete lateral: magnitudes calculadas con su
% incertidumbre.  Uso: \begin{destacado}[rótulo opcional] ... \end{destacado}
\newtcolorbox{destacado}[1][]{%
  enhanced, sharp corners, frame hidden, colback=azulfondo,
  borderline west={1.6pt}{0pt}{azulacero},
  left=10pt, right=8pt, top=2pt, bottom=4pt,
  before skip=8pt, after skip=10pt,
  before upper={\ifx\relax#1\relax\else{\scshape\small\color{azulacero}#1}\par\vspace{-4pt}\fi}}

\crefname{proposicion}{proposición}{proposiciones}
\Crefname{proposicion}{Proposición}{Proposiciones}
\crefname{teorema}{teorema}{teoremas}
\Crefname{teorema}{Teorema}{Teoremas}
\crefname{definicion}{definición}{definiciones}
\Crefname{definicion}{Definición}{Definiciones}
\crefname{resultado}{resultado}{resultados}
\Crefname{resultado}{Resultado}{Resultados}
\crefname{conclusion}{conclusión}{conclusiones}
\crefname{hipotesis}{hipótesis}{hipótesis}
\crefname{calculo}{cálculo}{cálculos}
\crefname{estimacion}{estimación}{estimaciones}
\crefname{observacion}{observación}{observaciones}
\Crefname{observacion}{Observación}{Observaciones}
\crefname{section}{sección}{secciones}
\Crefname{section}{Sección}{Secciones}
\crefname{subsection}{sección}{secciones}
\Crefname{subsection}{Sección}{Secciones}
\crefname{table}{tabla}{tablas}
\Crefname{table}{Tabla}{Tablas}
\crefname{figure}{figura}{figuras}
\Crefname{figure}{Figura}{Figuras}
\crefname{equation}{ecuación}{ecuaciones}
\Crefname{equation}{Ecuación}{Ecuaciones}
\crefformat{equation}{ecuación~(#2#1#3)}
\Crefformat{equation}{Ecuación~(#2#1#3)}
\crefrangeformat{equation}{ecuaciones~(#3#1#4)--(#5#2#6)}
\crefmultiformat{equation}{ecuaciones~(#2#1#3)}{ y~(#2#1#3)}{, (#2#1#3)}{ y~(#2#1#3)}

% --- Estilo de gráficas de datos ----------------------------------------------
\pgfplotsset{
  datos/.style={
    tick label style={font=\scriptsize}, label style={font=\small},
    axis lines=left, axis line style={grispizarra},
    grid=major, grid style={grisclaro, line width=0.4pt},
    every tick/.style={grispizarra},
    legend style={font=\scriptsize, draw=none, fill=white, fill opacity=0.9,
                  text opacity=1, cells={anchor=west}},
    legend cell align=left,
    /pgf/number format/use comma, /pgf/number format/1000 sep={\,},
    scaled ticks=false,
  },
  puntoA/.style={only marks, mark=*, mark size=2.1pt, azulprof,
    mark options={fill=azulprof, draw=white, line width=0.3pt},
    error bars/.cd, y dir=both, y explicit,
    error bar style={line width=0.6pt, azulprof}, error mark options={rotate=90, mark size=2.2pt, line width=0.6pt, azulprof}},
  puntoB/.style={only marks, mark=square*, mark size=1.9pt, azulacero,
    mark options={fill=azulclaro, draw=azulacero, line width=0.5pt},
    error bars/.cd, y dir=both, y explicit,
    error bar style={line width=0.6pt, azulacero}, error mark options={rotate=90, mark size=2.2pt, line width=0.6pt, azulacero}},
  ajusteB/.style={azulacero, line width=1.1pt, densely dashed},
  % leyenda en una fila sobre el gráfico, para que no tape los datos
  leyenda arriba/.style={legend columns=-1,
    legend style={at={(0.5,1.03)}, anchor=south,
      /tikz/every even column/.append style={column sep=9pt}}},
  ajuste/.style={azulprof, line width=1.1pt},
  teorica/.style={acento, line width=0.9pt, densely dashed},
}

% Referencia a un panel de figura: «figura 2\,(a)», enlazada en bloque
\newcommand{\figpanel}[2]{\hyperref[#1]{figura~\ref*{#1}\,\textup{(#2)}}}

% Pies de figura y tabla
\usepackage[font=small, labelfont={sc,color=azulprof}, labelsep=period,
            margin=0.6cm, justification=justified]{caption}
\usepackage{subcaption}
\captionsetup[sub]{font=small, labelfont={sc,color=azulprof}, labelformat=parens,
                   margin=0pt, justification=centering}

% =========================================================================
%                     VARIABLES PARA LA PORTADA
% =========================================================================
\newcommand{\Universidad}{Universidad de Salamanca}
\newcommand{\Facultad}{Facultad de Ciencias}
\newcommand{\Titulacion}{Grado en Física}
\newcommand{\Asignatura}{Laboratorio de Mecánica y Ondas}
\newcommand{\NumPractica}{Práctica 2}
\newcommand{\TituloPractica}{Ondas estacionarias}
\newcommand{\SubtituloPractica}{Modos normales de una cuerda tensa}
\newcommand{\tipotrabajo}{Informe de laboratorio}
\newcommand{\miNombre}{David Gómez García \\ Daniel Martín Guerrero}
\newcommand{\miTutor}{Mario Amado \\ Ivan de Martino \\ Manuel Pino}
\newcommand{\miUbicacion}{Salamanca}

% =========================================================================
%                     DATOS DEL ARTÍCULO (cabecera)
% =========================================================================
\title{Ondas estacionarias en una cuerda}
\author{David Gómez García}
\author{Daniel Martín Guerrero}
\address{Grado en Física, Facultad de Ciencias, Universidad de Salamanca,
         Salamanca, España}
\date{}

% =========================================================================
%                                DOCUMENTO
% =========================================================================
\begin{document}
% --- PORTADA ----------------------------------------------------------------
\pagenumbering{alph}
%%%%%%%%%%%%%%%%%%%%%%%%%%%%%%%%%%%%%%%%%%%%%%%%%%%%%%%%%%%%%%%%%%%%%%%%
% Portada — estilo artículo científico, motivo: onda estacionaria vista
% con el estroboscopio (instantáneas superpuestas del modo n = 3).
% Misma cabecera, bloque de título y pie que los informes anteriores.
%%%%%%%%%%%%%%%%%%%%%%%%%%%%%%%%%%%%%%%%%%%%%%%%%%%%%%%%%%%%%%%%%%%%%%%%

\begin{titlepage}
\thispagestyle{empty}
\pagecolor{white}

\begin{tikzpicture}[remember picture, overlay]
  \coordinate (O) at (current page.south west);

  % ----------------------------------------------------------------
  % Cabecera: logotipos y filete
  % ----------------------------------------------------------------
  \node[anchor=north west, inner sep=0] at ($(O)+(2.2cm,27.9cm)$)
    {\includegraphics[height=1.42cm]{figuras/logo-usal.png}};
  \node[anchor=north east, inner sep=0] at ($(O)+(18.8cm,27.9cm)$)
    {\includegraphics[height=1.42cm]{figuras/logo-ciencias.png}};
  \draw[grispizarra!60, line width=0.4pt]
    ($(O)+(2.2cm,26.0cm)$) -- ($(O)+(18.8cm,26.0cm)$);
  \node[anchor=north west, inner sep=0, text=grispizarra, font=\scshape\small]
    at ($(O)+(2.2cm,25.75cm)$)
    {\textls[80]{\Universidad\ \textbar\ \Facultad}};
  \node[anchor=north east, inner sep=0, text=grispizarra, font=\scshape\small]
    at ($(O)+(18.8cm,25.75cm)$)
    {\textls[80]{\Titulacion}};

  % ----------------------------------------------------------------
  % Motivo gráfico: modo n = 3 «congelado» por el estroboscopio.
  % Cada traza es y(x,t_j) = A sin(3 pi x / L) cos(phi_j); su conjunto
  % dibuja la envolvente, con los nodos inmóviles.
  % ----------------------------------------------------------------
  \begin{scope}[shift={($(O)+(2.7cm,19.4cm)$)}, x=1cm, y=1cm]
    \def\Lc{15.6}\def\Ac{2.35}
    % envolvente tenue
    \fill[azulclaro!22] plot[domain=0:\Lc, samples=200, smooth] (\x,{\Ac*sin(3*180*\x/\Lc)})
      -- plot[domain=\Lc:0, samples=200, smooth] (\x,{-\Ac*sin(3*180*\x/\Lc)}) -- cycle;
    % posición de equilibrio
    \draw[grispizarra!45, line width=0.4pt, dashed] (-0.5,0) -- (\Lc+0.5,0);
    % instantáneas
    \foreach \j in {0,...,12}{
      \pgfmathsetmacro{\c}{cos(15*\j)}
      \pgfmathsetmacro{\pc}{25+70*abs(\c)}
      \pgfmathsetmacro{\lw}{0.25+0.75*abs(\c)}
      \draw[azulprof!\pc!azulclaro, line width=\lw pt]
        plot[domain=0:\Lc, samples=160, smooth] (\x,{\Ac*\c*sin(3*180*\x/\Lc)});}
    % vientres
    \foreach \k in {1,3,5}{
      \draw[acento!80, line width=0.4pt, densely dotted]
        ({\Lc*\k/6},{-\Ac-0.15}) -- ({\Lc*\k/6},{\Ac+0.15});}
    % nodos
    \foreach \k in {0,...,3}{
      \fill[white] ({\Lc*\k/3},0) circle (3.4pt);
      \fill[azulprof] ({\Lc*\k/3},0) circle (2.4pt);}
    % extremos: excéntrica del motor y varilla de apoyo
    \draw[azulprof, line width=0.7pt, fill=white] (-0.55,0) circle (0.42);
    \fill[azulprof] (-0.55,0) circle (1.2pt);
    \fill[acento] (-0.13,0) circle (1.8pt);
    \draw[grispizarra, line width=1.6pt] (\Lc,-2.9) -- (\Lc,-0.12);
    % longitud de onda
    \draw[grispizarra, line width=0.45pt, |<->|]
      (0,{\Ac+0.75}) -- ({2*\Lc/3},{\Ac+0.75})
      node[midway, fill=white, inner sep=2pt, text=azulprof] {$\lambda_{3}=2L/3$};
    % distancia entre apoyos
    \draw[grispizarra, line width=0.45pt, |<->|]
      (0,{-\Ac-0.75}) -- (\Lc,{-\Ac-0.75})
      node[midway, fill=white, inner sep=2pt, text=azulprof] {$L$};
    % rótulos discretos, al modo de «tú estás aquí»
    \draw[grispizarra!70, line width=0.35pt] ({\Lc/3},0) -- ({\Lc/3+0.55},{-\Ac-0.25});
    \node[anchor=west, text=grispizarra, font=\itshape\small] at ({\Lc/3+0.6},{-\Ac-0.27}) {nodo};
    \draw[grispizarra!70, line width=0.35pt] ({\Lc*5/6},{-\Ac}) -- ({\Lc*5/6-1.2},{-\Ac+0.35});
    \fill[acento] ({\Lc*5/6},{-\Ac}) circle (1.3pt);
    \node[anchor=east, text=grispizarra, font=\itshape\small] at ({\Lc*5/6-1.25},{-\Ac+0.38}) {vientre};
    \node[anchor=north, text=grispizarra, font=\itshape\small] at (-0.55,-0.55) {motor};
  \end{scope}

  % ----------------------------------------------------------------
  % Fila de modos normales (como la línea del tiempo de Cosmosofía)
  % ----------------------------------------------------------------
  \begin{scope}[shift={($(O)+(0,13.55cm)$)}, x=1cm, y=1cm]
    \draw[grispizarra!60, line width=0.4pt] (2.2,0) -- (18.8,0);
    \foreach \n in {1,...,5}{
      \pgfmathsetmacro{\xa}{2.55+3.27*(\n-1)}
      \pgfmathsetmacro{\pc}{35+15*\n}
      \fill[azulclaro!40] plot[domain=0:2.6, samples=60, smooth] ({\xa+\x},{0.5*sin(\n*180*\x/2.6)})
        -- plot[domain=2.6:0, samples=60, smooth] ({\xa+\x},{-0.5*sin(\n*180*\x/2.6)}) -- cycle;
      \draw[azulprof!\pc!azulclaro, line width=0.9pt]
        plot[domain=0:2.6, samples=60, smooth] ({\xa+\x},{0.5*sin(\n*180*\x/2.6)});
      \draw[azulprof!\pc!azulclaro, line width=0.5pt, densely dashed]
        plot[domain=0:2.6, samples=60, smooth] ({\xa+\x},{-0.5*sin(\n*180*\x/2.6)});
      \fill[azulprof] (\xa,0) circle (1.2pt);
      \fill[azulprof] ({\xa+2.6},0) circle (1.2pt);
      \node[anchor=north, text=grispizarra, font=\footnotesize] at ({\xa+1.3},-0.62) {$n=\n$};}
    \node[anchor=north, text=grispizarra, font=\scshape\scriptsize] at (10.5,-1.08)
      {\textls[160]{modos normales de una cuerda fija por sus extremos}};
  \end{scope}

  % ----------------------------------------------------------------
  % Bloque del título
  % ----------------------------------------------------------------
  \node[anchor=north west, inner sep=0, text width=16.6cm, align=left]
    at ($(O)+(2.2cm,11.1cm)$)
    {{\color{azulacero}\scshape\large\textls[120]{\Asignatura\ \textbar\ \NumPractica}}\\[12pt]
     {\color{azulprof}\fontsize{33}{38}\selectfont\bfseries\spaceskip=0.28em\relax\TituloPractica}\\[10pt]
     {\color{azulprof}\itshape\Large\SubtituloPractica}\\[12pt]
     {\color{azulacero}\rule{2.5cm}{1.4pt}}\\[9pt]
     {\color{grispizarra}\large
      $\lambda_{n}=\dfrac{2L}{n}$
      \quad\textperiodcentered\quad
      $\lambda^{2}=\dfrac{T}{\mu f^{2}}$}};

  % ----------------------------------------------------------------
  % Pie: tipo de trabajo, autores, tutores, lugar y fecha
  % ----------------------------------------------------------------
  \draw[grispizarra!60, line width=0.4pt]
    ($(O)+(2.2cm,4.6cm)$) -- ($(O)+(18.8cm,4.6cm)$);

  \node[anchor=north west, inner sep=0, text width=4.6cm, align=left]
    at ($(O)+(2.2cm,4.25cm)$)
    {{\scshape\small\color{azulacero}\textls[80]{Documento}}\\[3pt]
     {\large\bfseries\color{azulprof}\tipotrabajo}};

  \node[anchor=north west, inner sep=0, text width=5.6cm, align=left]
    at ($(O)+(7.2cm,4.25cm)$)
    {{\scshape\small\color{azulacero}\textls[80]{Autores}}\\[3pt]
     {\large\color{black!85}\miNombre}};

  \node[anchor=north west, inner sep=0, text width=5cm, align=left]
    at ($(O)+(13.8cm,4.25cm)$)
    {{\scshape\small\color{azulacero}\textls[80]{Tutores}}\\[3pt]
     {\large\color{black!85}\miTutor}};

  \node[anchor=south east, inner sep=0, text=grispizarra, font=\scshape\small]
    at ($(O)+(18.8cm,1.2cm)$)
    {\textls[100]{\miUbicacion, \number\year}};
  \fill[azulprof] ($(O)+(0,0)$) rectangle ($(O)+(21cm,0.45cm)$);
  \fill[azulacero] ($(O)+(0,0)$) rectangle ($(O)+(7cm,0.45cm)$);
  \fill[azulclaro] ($(O)+(0,0)$) rectangle ($(O)+(3cm,0.45cm)$);
\end{tikzpicture}

\end{titlepage}

\clearpage
\pagenumbering{arabic}
\setcounter{page}{1}

\begin{abstract}
Excitamos con un motor una cuerda tensa sujeta por sus dos extremos y
buscamos, variando la tensión, las condiciones en las que se forman ondas
estacionarias con uno, dos o más vientres. Para cada modo la longitud de onda
queda fijada por la geometría, $\lambda_n=2L/n$, mientras que la frecuencia,
la del motor, se mide con un estroboscopio. Representando $\lambda^{2}$ frente
a la tensión $T$ para dos distancias entre apoyos y dos cuerdas, una simple y
otra doble, comprobamos que la relación es lineal, que la pendiente
$1/(\mu f^{2})$ no depende de la distancia y que se reduce a la mitad al
doblar la densidad lineal. De las pendientes estimamos la densidad lineal de
la cuerda y la comparamos con la que obtenemos pesándola. Los resultados se
expresan con su incertidumbre típica y ampliada al \SI{95}{\percent},
siguiendo \emph{Notes on Uncertainty in Measurements}.
\end{abstract}

\maketitle
\tableofcontents

% =========================================================================
\newpage
\seccion{Introducción}{La música de una cuerda tensa}\label{sec:introduccion}
% =========================================================================
Cuando pulsamos la cuerda de una guitarra no suena cualquier frecuencia: solo
aquellas para las que la onda que viaja por la cuerda, al reflejarse en los
extremos, se superpone consigo misma de forma que los extremos permanezcan
quietos. El resultado no es una onda que avanza, sino una \emph{onda
estacionaria}, con puntos que no se mueven nunca (los nodos) y otros que
oscilan con amplitud máxima (los vientres o antinodos)~\cite{FV5-ECJ,FRENCH}.

En esta práctica hacemos el experimento al revés: fijamos la frecuencia, la de
un motor, y variamos la tensión de la cuerda hasta que la onda «encaja» entre
los dos apoyos con uno, dos o más vientres. Como la velocidad de propagación
depende de la tensión y de la densidad lineal de la cuerda, la tensión a la
que aparece cada modo nos permite medir esa densidad sin pesar la
cuerda~\cite{GUION-P2}.

\subsection{Objetivos}\label{sec:objetivos}
\begin{itemize}
  \item Determinar la longitud de onda de los modos de vibración de una cuerda
        en función de la tensión, de la densidad lineal y de la distancia
        entre los puntos de apoyo.
  \item Comprobar que $\lambda^{2}$ crece linealmente con la tensión y
        comparar las pendientes para distintas distancias y distintas
        densidades lineales.
  \item Estimar la densidad lineal de la cuerda a partir de las pendientes y
        compararla con la obtenida con la báscula.
\end{itemize}
Las tres cuestiones se resuelven con la \cref{eq:lineal}, que relaciona
linealmente $\lambda^{2}$ con $T$.

\subsection{Material empleado}\label{sec:material}
Hemos utilizado el siguiente instrumental~\cite{GUION-P2}:
\begin{itemize}
  \item Motor generador de ondas, con un dispositivo excéntrico.
  \item Cuerda (simple y doble) y varilla vertical de apoyo.
  \item Dinamómetro sobre un soporte desplazable.
  \item Estroboscopio.
  \item Regla milimetrada y báscula.
\end{itemize}

\subsection{Montaje experimental}\label{sec:montaje}
Un extremo de la cuerda se ata a la excéntrica del motor, que la hace
oscilar, y el otro, tras pasar por una varilla vertical, al dinamómetro
(\cref{fig:montaje}). La distancia $L$ entre la excéntrica y la varilla es la
longitud de cuerda que vibra, y la varilla puede fijarse en dos posiciones
del banco, $L_1>L_2$. Desplazando lentamente el soporte del dinamómetro
variamos la tensión hasta que aparece cada modo estacionario, y con el
estroboscopio medimos la frecuencia de vibración: cuando la frecuencia de los
destellos coincide con la del motor, la cuerda parece congelada.

\begin{figure}[ht]
\centering
\begin{tikzpicture}[>=Stealth, x=1cm, y=1cm, font=\small]
  % -------- banco --------
  \fill[grisclaro] (-2.2,-1.55) rectangle (12.6,-1.3);
  \draw[grispizarra, line width=0.5pt] (-2.2,-1.55) rectangle (12.6,-1.3);
  \foreach \x in {-1.8,-1.2,...,12.3}{\draw[grispizarra!50, line width=0.3pt] (\x,-1.55)--(\x,-1.3);}
  % -------- motor con excéntrica --------
  \draw[azulprof, fill=azulclaro!45, line width=0.7pt, rounded corners=2pt] (-1.95,-1.3) rectangle (-0.45,0.95);
  \node[azulprof, font=\footnotesize] at (-1.2,-0.8) {motor};
  \draw[azulprof, fill=white, line width=0.7pt] (-0.85,0.25) circle (0.38);
  \fill[azulprof] (-0.85,0.25) circle (0.05);
  \fill[acento] (-0.66,0.25) circle (0.07);
  \draw[acento, line width=0.5pt, ->] (-0.85,0.25) ++(35:0.55) arc[start angle=35, end angle=150, radius=0.55];
  \node[grispizarra, font=\scriptsize\itshape, anchor=south] at (-1.2,1.0) {excéntrica};
  % -------- varilla de apoyo (L1) y posición alternativa (L2) --------
  \draw[grispizarra, line width=1.6pt] (8.6,-1.3) -- (8.6,0.15);
  \fill[grispizarra] (8.6,0.25) circle (0.1);
  \draw[grispizarra!50, line width=1.2pt, densely dashed] (6.4,-1.3) -- (6.4,0.15);
  \node[grispizarra, font=\scriptsize\itshape, anchor=north] at (6.4,-1.62) {posición $L_2$};
  \node[grispizarra, font=\scriptsize\itshape, anchor=north] at (8.75,-1.62) {varilla (posición $L_1$)};
  % -------- onda estacionaria n = 3 --------
  \def\xa{-0.66}\def\xb{8.6}\def\yc{0.25}\def\amp{0.62}
  \fill[azulclaro!35] plot[domain=0:1, samples=120, smooth] ({\xa+(\xb-\xa)*\x},{\yc+\amp*sin(3*180*\x)})
    -- plot[domain=1:0, samples=120, smooth] ({\xa+(\xb-\xa)*\x},{\yc-\amp*sin(3*180*\x)}) -- cycle;
  \draw[azulacero, line width=0.6pt, densely dashed] plot[domain=0:1, samples=120, smooth]
      ({\xa+(\xb-\xa)*\x},{\yc-\amp*sin(3*180*\x)});
  \draw[azulprof, line width=1.1pt] plot[domain=0:1, samples=120, smooth]
      ({\xa+(\xb-\xa)*\x},{\yc+\amp*sin(3*180*\x)});
  \foreach \k in {0,1,2,3}{\fill[azulprof] ({\xa+(\xb-\xa)*\k/3},\yc) circle (1.6pt);}
  \foreach \k in {1,3,5}{\draw[acento, line width=0.4pt, densely dotted]
      ({\xa+(\xb-\xa)*\k/6},{\yc-\amp}) -- ({\xa+(\xb-\xa)*\k/6},{\yc+\amp});}
  \draw[grispizarra!70, line width=0.35pt] ({\xa+(\xb-\xa)/3},\yc) -- ({\xa+(\xb-\xa)/3+0.35},{\yc-\amp-0.12});
  \node[azulprof, font=\scriptsize, anchor=north west] at ({\xa+(\xb-\xa)/3+0.3},{\yc-\amp-0.08}) {nodo};
  \node[acento, font=\scriptsize, anchor=south] at ({\xa+(\xb-\xa)/2},{\yc+\amp+0.02}) {vientre};
  % -------- cuerda hasta el dinamómetro --------
  \draw[azulprof, line width=1.1pt] (8.6,0.35) -- (10.0,0.35);
  \draw[grispizarra, fill=white, line width=0.7pt, rounded corners=1.5pt] (10.0,0.12) rectangle (11.75,0.58);
  \foreach \x in {10.2,10.4,...,11.6}{\draw[grispizarra!70, line width=0.3pt] (\x,0.12) -- (\x,0.27);}
  \node[grispizarra, font=\scriptsize\itshape, anchor=south] at (10.88,0.62) {dinamómetro, $T$};
  \draw[grispizarra, line width=1.4pt] (12.0,-1.3) -- (12.0,0.35) -- (11.75,0.35);
  \draw[->, acento, line width=0.8pt] (12.2,-0.35) -- (12.85,-0.35);
  \node[acento, font=\scriptsize\itshape, align=left, anchor=west] at (12.1,-0.85) {soporte\\móvil};
  % -------- cota L --------
  \draw[grispizarra, line width=0.45pt, |<->|] (\xa,-0.95) -- (\xb,-0.95)
    node[midway, fill=white, inner sep=2pt, text=black!80] {$L$};
  % -------- estroboscopio --------
  \draw[azulprof, fill=azulclaro!45, line width=0.7pt, rounded corners=2pt] (4.9,1.75) rectangle (7.7,2.5);
  \node[azulprof, font=\scriptsize] at (6.3,2.12) {estroboscopio, $f$};
  \foreach \ang in {-125,-105,-85}{\draw[acento!80, line width=0.4pt] (6.3,1.75) -- ++(\ang:0.7);}
\end{tikzpicture}
\caption{Esquema del montaje, con la cuerda vibrando en el modo $n=3$. La
excéntrica del motor hace oscilar la cuerda, que pasa por la varilla vertical
y llega al dinamómetro, sujeto a un soporte que se desplaza para variar la
tensión $T$. La longitud que vibra, $L$, va de la excéntrica a la varilla, que
puede ocupar dos posiciones ($L_1>L_2$). Esquema sin escala.}
\label{fig:montaje}
\end{figure}

\ornamento

% =========================================================================
\newpage
\seccion{Fundamentos teóricos}{De la ecuación de ondas a los modos normales}\label{sec:fundamentos}
% =========================================================================
\subsection{Ondas transversales en una cuerda tensa}\label{sec:ecuacion}
Consideremos una cuerda homogénea, de densidad lineal $\mu$, sometida a una
tensión $T$ y desplazada ligeramente de su posición de equilibrio. Llamamos
$y(x,t)$ al desplazamiento transversal del punto de abscisa $x$ en el instante
$t$. Un elemento de longitud $\dif x$ está sometido a las tensiones de sus dos
extremos, tangentes a la cuerda (\cref{fig:elemento}); para desplazamientos
pequeños su módulo apenas cambia y sus componentes horizontales se
compensan, pero no las verticales.

\begin{figure}[ht]
\centering
\begin{tikzpicture}[>=Stealth, x=1cm, y=1cm, font=\small]
  % cuerda deformada: y = 1,6 sin(30 x) (grados); elemento entre xA y xB
  \def\xA{0.9}\def\xB{4.3}\def\Am{1.6}
  \pgfmathsetmacro{\yA}{\Am*sin(30*\xA)}
  \pgfmathsetmacro{\yB}{\Am*sin(30*\xB)}
  \pgfmathsetmacro{\aA}{atan(\Am*30*pi/180*cos(30*\xA))}
  \pgfmathsetmacro{\aB}{atan(\Am*30*pi/180*cos(30*\xB))}
  \draw[azulclaro, line width=0.7pt, domain=-0.6:7.2, samples=120, smooth] plot (\x,{\Am*sin(30*\x)});
  \draw[azulprof, line width=1.8pt, domain=\xA:\xB, samples=60, smooth] plot (\x,{\Am*sin(30*\x)});
  \draw[grispizarra!45, line width=0.4pt, dashed] (-0.6,0) -- (7.2,0);
  \coordinate (A) at (\xA,\yA);
  \coordinate (B) at (\xB,\yB);
  % tensiones, tangentes a la cuerda
  \draw[->, acento, line width=1.1pt] (A) -- ++({180+\aA}:1.7) node[anchor=north east, text=acento, inner sep=1pt] {$T$};
  \draw[->, acento, line width=1.1pt] (B) -- ++(\aB:1.7) node[anchor=north west, text=acento, inner sep=1pt] {$T$};
  % ángulos con la horizontal
  \draw[grispizarra!70, line width=0.4pt, densely dotted] (A) -- ++(-1.55,0);
  \draw[grispizarra, line width=0.6pt] (A) ++(180:1.05) arc[start angle=180, end angle={180+\aA}, radius=1.05];
  \node[grispizarra, anchor=east, font=\footnotesize] at ($(A)+(-1.6,-0.05)$) {$\theta(x)$};
  \draw[grispizarra!70, line width=0.4pt, densely dotted] (B) -- ++(1.55,0);
  \draw[grispizarra, line width=0.6pt] (B) ++(0:1.05) arc[start angle=0, end angle=\aB, radius=1.05];
  \node[grispizarra, anchor=west, font=\footnotesize] at ($(B)+(1.6,0.05)$) {$\theta(x+\dif x)$};
  \fill[azulprof] (A) circle (2pt);
  \fill[azulprof] (B) circle (2pt);
  \node[azulprof, anchor=south] at (3.0,{\Am+0.12}) {$\dif m=\mu\,\dif x$};
  % cotas
  \draw[grispizarra, densely dotted, line width=0.4pt] (A) -- (\xA,-2.05);
  \draw[grispizarra, densely dotted, line width=0.4pt] (B) -- (\xB,-2.05);
  \draw[grispizarra, line width=0.45pt, |<->|] (\xA,-2.05) -- (\xB,-2.05)
    node[midway, fill=white, inner sep=2pt, text=black!80] {$\dif x$};
  \draw[->, grispizarra, line width=0.5pt] (-1.9,-2.6) -- (7.5,-2.6) node[right, text=black!80] {$x$};
  \draw[->, grispizarra, line width=0.5pt] (-1.9,-2.6) -- (-1.9,2.3) node[above, text=black!80] {$y$};
\end{tikzpicture}
\caption{Elemento de cuerda de longitud $\dif x$ y masa $\mu\,\dif x$. Las
tensiones en sus extremos son tangentes a la cuerda y forman ángulos
$\theta(x)$ y $\theta(x+\dif x)$ con la horizontal; su resultante vertical es
la que acelera el elemento. Desplazamientos exagerados.}
\label{fig:elemento}
\end{figure}

Aplicando la segunda ley de Newton a la componente vertical, y usando que para
ángulos pequeños $\sin\theta\simeq\tan\theta=\partial y/\partial x$,
\begin{equation}\label{eq:newton}
  \mu\,\dif x\,\pdd{y}{t}
  =T\bigl[\sin\theta(x+\dif x)-\sin\theta(x)\bigr]
  \simeq T\left[\pd{y}{x}\Big|_{x+\dif x}-\pd{y}{x}\Big|_{x}\right]
  =T\,\pdd{y}{x}\,\dif x ,
\end{equation}
es decir, la ecuación de ondas~\cite{GUION-P2,FV6-ECJ}
\begin{equation}\label{eq:ondas}
  \pdd{y}{t}=\frac{T}{\mu}\,\pdd{y}{x}
  \qquad\Longrightarrow\qquad
  v=\sqrt{\frac{T}{\mu}} ,
\end{equation}
donde la comparación con la forma general
$\partial^{2}y/\partial t^{2}=v^{2}\,\partial^{2}y/\partial x^{2}$ identifica
la velocidad de propagación $v$. Tiene sentido: una cuerda más tensa responde
con más fuerza a una deformación, y una cuerda más pesada, con más inercia.

Para una onda armónica de frecuencia $f$, la longitud de onda es la distancia
que recorre la perturbación en un periodo, $\lambda=v/f$, así que
\begin{equation}\label{eq:lambda}
  \lambda=\frac{1}{f}\sqrt{\frac{T}{\mu}} .
\end{equation}

\subsection{Ondas estacionarias}\label{sec:estacionarias}
Una onda que viaja por la cuerda y se refleja en un extremo fijo da lugar a
una onda reflejada de igual amplitud y longitud de onda que se propaga en
sentido contrario. Su superposición es~\cite{GUION-P2,FV5-ECJ}
\begin{equation}\label{eq:superposicion}
  y(x,t)=A\sin(kx-\omega t)+A\sin(kx+\omega t)=2A\sin(kx)\cos(\omega t),
\end{equation}
con $k=2\pi/\lambda$ y $\omega=2\pi f$. La dependencia espacial y la temporal
quedan separadas: todos los puntos oscilan en fase, cada uno con amplitud
$\lvert2A\sin(kx)\rvert$ (\cref{fig:estacionaria}). Los puntos con
$\sin(kx)=0$ no se mueven nunca, mientras que en los que cumplen
$\lvert\sin(kx)\rvert=1$ la amplitud es máxima.

\begin{referencia}[Nodos, vientres y cuantización]{GUION-P2}\label{ref:modos}
Se denominan antinodos a los puntos de la onda estacionaria que tienen amplitud
máxima [\dots] Por lo tanto, se puede decir que los modos de vibración están
cuantizados, y son múltiplos del modo fundamental $k_1=\pi/L$, es decir,
$k_n=n\,k_1$.
\end{referencia}

\begin{figure}[ht]
\centering
\begin{tikzpicture}[>=Stealth, x=1cm, y=1cm, font=\small]
  \def\Lw{12}\def\Aw{1.25}
  \fill[azulclaro!25] plot[domain=0:\Lw, samples=160, smooth] (\x,{\Aw*sin(2*180*\x/\Lw)})
    -- plot[domain=\Lw:0, samples=160, smooth] (\x,{-\Aw*sin(2*180*\x/\Lw)}) -- cycle;
  \draw[grispizarra!45, line width=0.4pt, dashed] (-0.3,0) -- (\Lw+0.3,0);
  \foreach \j in {0,...,8}{
    \pgfmathsetmacro{\c}{cos(22.5*\j)}
    \pgfmathsetmacro{\pc}{25+70*abs(\c)}
    \draw[azulprof!\pc!azulclaro, line width=0.7pt] plot[domain=0:\Lw, samples=120, smooth] (\x,{\Aw*\c*sin(2*180*\x/\Lw)});}
  \foreach \k in {0,1,2}{
    \fill[white] ({\Lw*\k/2},0) circle (3pt); \fill[azulprof] ({\Lw*\k/2},0) circle (2.2pt);
    \node[azulprof, anchor=north, font=\footnotesize] at ({\Lw*\k/2},-0.12) {N};}
  \foreach \k in {1,3}{
    \draw[acento, line width=0.4pt, densely dotted] ({\Lw*\k/4},{-\Aw-0.1}) -- ({\Lw*\k/4},{\Aw+0.1});
    \node[acento, anchor=south, font=\footnotesize] at ({\Lw*\k/4},{\Aw+0.1}) {A};}
  \draw[grispizarra, line width=0.45pt, |<->|] (0,{-\Aw-0.55}) -- ({\Lw/2},{-\Aw-0.55})
    node[midway, fill=white, inner sep=2pt, text=black!80] {$\lambda/2$};
  \draw[grispizarra, line width=0.45pt, |<->|] ({\Lw/4},{\Aw+0.75}) -- ({3*\Lw/4},{\Aw+0.75})
    node[midway, fill=white, inner sep=2pt, text=black!80] {$\lambda/2$};
  \draw[grispizarra, line width=0.45pt, <->] ({\Lw*3/4+0.35},{-\Aw}) -- ({\Lw*3/4+0.35},{\Aw})
    node[midway, right, text=black!80] {$4A$};
\end{tikzpicture}
\caption{Onda estacionaria $y=2A\sin(kx)\cos(\omega t)$: instantáneas en
distintos momentos de un periodo, como las que congela el estroboscopio. Los
nodos (N) permanecen quietos y los vientres (A) oscilan con amplitud $2A$; dos
nodos, o dos vientres, consecutivos distan $\lambda/2$.}
\label{fig:estacionaria}
\end{figure}

\subsection{Modos de una cuerda con los extremos fijos}\label{sec:modos}
Si la cuerda está sujeta en $x=0$ y $x=L$, ambos extremos han de ser nodos:
$\sin(kL)=0$. Solo son posibles, por tanto, los números de onda
$k_n=n\pi/L$ o, equivalentemente,
\begin{equation}\label{eq:modos}
  \lambda_n=\frac{2L}{n},\qquad n=1,2,3,\dots
\end{equation}
donde $n$ es el número de vientres (\cref{fig:modos}). Es la cuantización de
la que habla la \cref{ref:modos}: entre los apoyos ha de caber un número
entero de semilongitudes de onda~\cite{GUION-P2,FRENCH}.

\begin{figure}[ht]
\centering
\begin{tikzpicture}[x=1cm, y=1cm, font=\small]
  \foreach \n/\y in {1/0,2/-1.1,3/-2.2,4/-3.3,5/-4.4}{
    \fill[azulclaro!35] plot[domain=0:1, samples=100, smooth] ({7*\x},{\y+0.36*sin(\n*180*\x)})
      -- plot[domain=1:0, samples=100, smooth] ({7*\x},{\y-0.36*sin(\n*180*\x)}) -- cycle;
    \draw[azulprof, line width=0.9pt] plot[domain=0:1, samples=100, smooth] ({7*\x},{\y+0.36*sin(\n*180*\x)});
    \draw[azulacero, line width=0.5pt, densely dashed] plot[domain=0:1, samples=100, smooth] ({7*\x},{\y-0.36*sin(\n*180*\x)});
    \fill[grispizarra] (0,\y) circle (1.5pt); \fill[grispizarra] (7,\y) circle (1.5pt);
    \node[anchor=east, text=azulprof] at (-0.3,\y) {$n=\n$};
    \pgfmathsetmacro{\Tn}{1/(\n*\n)}
    \node[anchor=west, text=black!80] at (7.4,\y) {$\lambda_{\n}=\ifnum\n=1 2L\else\ifnum\n=2 L\else 2L/\n\fi\fi$};
    \fill[acento!70] (9.9,\y-0.11) rectangle ({9.9+2.4*\Tn},\y+0.11);
  }
  \node[anchor=south west, text=grispizarra, font=\scriptsize\itshape] at (9.9,0.3) {tensión necesaria, $T_n\propto 1/n^{2}$};
\end{tikzpicture}
\caption{Primeros modos de una cuerda de longitud $L$ con los extremos fijos.
A la derecha, la tensión que requiere cada modo a frecuencia fija,
$T_n=\mu f^{2}(2L/n)^{2}$, en proporción a la del fundamental: el modo $n=2$
necesita la cuarta parte de tensión que el $n=1$.}
\label{fig:modos}
\end{figure}

\subsection{Linealización}\label{sec:linealizacion}
En nuestro montaje la frecuencia es fija, la del motor, y lo que variamos es
la tensión. Elevando al cuadrado la \cref{eq:lambda}:

\begin{resultado}[Relación entre longitud de onda y tensión]\label{res:lineal}
Para una cuerda de densidad lineal $\mu$ que vibra con frecuencia $f$, el
cuadrado de la longitud de onda es proporcional a la tensión,
\begin{equation}\label{eq:lineal}
  \lambda^{2}=\frac{1}{\mu f^{2}}\,T
  \qquad\Longleftrightarrow\qquad
  \mu=\frac{1}{a\,f^{2}},
\end{equation}
donde $a=1/(\mu f^{2})$ es la pendiente de la recta $\lambda^{2}$ frente a $T$.
\end{resultado}

Con la \cref{eq:modos}, cada modo observado aporta un punto
$(T_n,\lambda_n^{2})$ de esa recta. De la \cref{eq:lineal} se siguen las dos
comparaciones que pide el guion~\cite{GUION-P2}:

\begin{observacion}[Qué esperar de las pendientes]\label{obs:pendientes}
(i)~Para la misma distancia $L$, la cuerda doble, de densidad $2\mu$, debe
dar una pendiente que sea la \emph{mitad} que la de la cuerda simple.
(ii)~Para la misma cuerda, la pendiente \emph{no depende} de $L$: cambiar la
distancia entre apoyos solo desplaza los puntos a lo largo de la misma recta,
porque modifica las longitudes de onda permitidas y, con ellas, las tensiones a
las que aparecen los modos.
\end{observacion}

\ornamento

% =========================================================================
\newpage
\seccion{Procedimiento experimental}{Cómo medimos}\label{sec:procedimiento}
% =========================================================================
Seguimos la metodología del guion~\cite{GUION-P2}:
\begin{enumerate}
  \item \emph{Preparación.} Fijamos la varilla en la posición $L_1$, atamos la
        cuerda simple a la excéntrica y al dinamómetro, y comprobamos que este
        marcaba cero sin tensión.
  \item \emph{Frecuencia.} Con el motor en marcha, medimos su frecuencia con
        el estroboscopio, buscando la frecuencia de destellos a la que la
        cuerda aparecía inmóvil.
  \item \emph{Barrido en tensión.} Desplazando lentamente el soporte del
        dinamómetro, buscamos las tensiones a las que se formaban ondas
        estacionarias con $n=1,2,3,\dots$ vientres, y anotamos cada una.
  \item \emph{Resto de configuraciones.} Repetimos el barrido con la cuerda
        doble en $L_1$, y con ambas cuerdas tras pasar la varilla a la
        posición $L_2$.
  \item \emph{Densidad.} Pesamos un trozo de cuerda de longitud conocida en la
        báscula.
\end{enumerate}

\begin{observacion}[Sobre las lecturas]\label{obs:lecturas}
El estroboscopio también deja la cuerda inmóvil a frecuencias submúltiplos de
la del motor, $f/2$, $f/3$\ldots, por lo que tomamos la más alta a la que se
veía una sola imagen quieta. Por otra parte, la tensión de los modos altos es
muy pequeña (del orden de centésimas de newton) y queda cerca del límite de
lo que el dinamómetro permite leer con fiabilidad; esas medidas son las más
delicadas de la práctica.
\end{observacion}

\subsection{Tratamiento de incertidumbres}\label{sec:incertidumbres}
Seguimos la \emph{Guide to the expression of uncertainty in
measurement}~\cite{GUM}, tal como se recoge en las notas del
laboratorio~\cite{NUM-ARG}. Para una magnitud $y=f(x_1,\dots,x_N)$ con
entradas no correlacionadas, la incertidumbre típica se propaga como
\begin{equation}\label{eq:propagacion}
  u^{2}(y)=\sum_{k=1}^{N}\left(\frac{\partial f}{\partial x_k}\right)^{\!2}u^{2}(x_k),
\end{equation}
y la incertidumbre ampliada al \SI{95}{\percent} es $\ext{y}=\Tcoef\,u(y)$,
con $\Tcoef$ el factor de cobertura de Student para $\nu$ grados de
libertad~\cite{NUM-ARG}. Cuando una incertidumbre procede solo de la
resolución $\Delta$ de un instrumento, la tomamos como
$u=\Delta/\sqrt{12}$, con infinitos grados de libertad. Las incertidumbres se
dan con una sola cifra significativa, salvo cuando la primera es un 1, en
cuyo caso se conservan dos; el valor se redondea a la misma posición
decimal~\cite{ERR-MP}, y en notación científica valor e incertidumbre
comparten la misma potencia de diez.

\begin{metodo}[Regresión lineal por mínimos cuadrados]\label{met:regresion}
La \cref{eq:lineal} tiene la forma $y=a\,x+b$, con $y=\lambda^{2}$ y $x=T$.
Para los $N$ pares $(x_i,y_i)$ de cada serie ajustamos la recta por mínimos
cuadrados ordinarios, sin ponderar~\cite{NUM-ARG,HUGHES}. Con
$S_{xx}=\sum_i(x_i-\mean{x})^{2}$ y $S_{xy}=\sum_i(x_i-\mean{x})(y_i-\mean{y})$,
\begin{equation}\label{eq:regresion}
  a=\frac{S_{xy}}{S_{xx}},
  \qquad
  b=\mean{y}-a\,\mean{x},
  \qquad
  s^{2}=\frac{1}{N-2}\sum_{i=1}^{N}\bigl(y_i-a\,x_i-b\bigr)^{2},
\end{equation}
y las incertidumbres típicas de los parámetros son
\begin{equation}\label{eq:uregresion}
  u(a)=\frac{s}{\sqrt{S_{xx}}},
  \qquad
  u(b)=u(a)\,\sqrt{\mean{x^{2}}},
\end{equation}
con $\nu=N-2$ grados de libertad. La ordenada $b$ debería ser compatible con
cero, y la calidad del ajuste la juzgamos con el coeficiente de correlación
$r$ y con los residuos. Los ajustes se han realizado con un pequeño programa
en Python.
\end{metodo}

\begin{calculo}[Longitud de onda]\label{calc:lambda}
Como $\lambda_n=2L/n$ con $n$ exacto, la única fuente de incertidumbre es la
distancia entre apoyos:
\begin{equation}\label{eq:ulambda}
  \frac{u(\lambda_n^{2})}{\lambda_n^{2}}=2\,\frac{u(L)}{L}.
\end{equation}
\end{calculo}

\begin{calculo}[Densidad lineal]\label{calc:mu}
Aplicando la \cref{eq:propagacion} a $\mu=1/(af^{2})$ y a $\mub=m/\ell$,
\begin{equation}\label{eq:umu}
  \frac{u(\mus)}{\mus}=\sqrt{\left(\frac{u(a)}{a}\right)^{\!2}+\left(\frac{2\,u(f)}{f}\right)^{\!2}},
  \qquad
  \frac{u(\mub)}{\mub}=\sqrt{\left(\frac{u(m)}{m}\right)^{\!2}+\left(\frac{u(\ell)}{\ell}\right)^{\!2}} .
\end{equation}
Para $\mus$, el número efectivo de grados de libertad se calcula con la
fórmula de Welch--Sat\-ter\-thwaite, teniendo en cuenta que la contribución de
$f$ procede de una resolución, con $\nu\to\infty$~\cite{NUM-ARG}.
\end{calculo}

% =========================================================================
\seccion{Análisis de datos y resultados}{Lo que dicen las medidas}\label{sec:analisis}
% =========================================================================
\subsection{Magnitudes de partida}\label{sec:partida}
La frecuencia de los destellos con la que la cuerda aparecía congelada fue de
\SI{43}{\hertz} en casi todas las configuraciones (y \SI{44}{\hertz} en
algunos de los modos más altos, dentro de la resolución de \SI{1}{\hertz}
del estroboscopio), de modo que tomamos $f=\SI{43,0 \pm 0,3}{\hertz}$.

De cada distancia entre apoyos anotamos dos lecturas: \SI{38,8}{} y
\SI{37,6}{\centi\meter} para $L_2$, y \SI{51,9}{} y \SI{49,8}{\centi\meter}
para $L_1$. La diferencia, mucho mayor que la resolución de la regla, refleja
que los extremos de la longitud que vibra no están perfectamente definidos,
sobre todo en la excéntrica, que oscila. Tomamos la primera lectura como valor
de $L$ y la semidiferencia como su incertidumbre típica
(\cref{tab:partida}).

Por último, la báscula dio $m=\SI{0,06}{\gram}$ para un trozo de cuerda de
$\ell=\SI{49}{\centi\meter}$, con resoluciones de \SI{0,01}{\gram} y
\SI{0,1}{\centi\meter}. Por la \cref{eq:umu}, la densidad lineal medida
directamente es
\begin{destacado}[Densidad lineal con la báscula]
\vspace{-6pt}
\begin{equation}\label{eq:mub}
  \mub=\frac{m}{\ell}=\num{1,22 \pm 0,06 e-4}\ \si{\kilogram\per\meter}
  \quad(\text{incertidumbre típica}),
\end{equation}
{\small con una incertidumbre ampliada al \SI{95}{\percent} de
$\num{0,12e-4}\ \si{\kilogram\per\meter}$, dominada por la resolución de la
pesada (\SI{0,01}{\gram} sobre \SI{0,06}{\gram}).}
\end{destacado}

\begin{table}[ht]
\centering
\caption{Magnitudes de partida con sus incertidumbres típicas.}
\label{tab:partida}
\renewcommand{\arraystretch}{1.2}
\begin{tabular}{@{}l l l@{}}
\toprule
Magnitud & Valor & Origen de la incertidumbre \\
\midrule
Frecuencia, $f$             & \SI{43,0 \pm 0,3}{\hertz}         & resolución del estroboscopio \\
Distancia corta, $L_2$      & \SI{38,8 \pm 0,6}{\centi\meter}   & dispersión de dos lecturas \\
Distancia larga, $L_1$      & \SI{51,9 \pm 1,0}{\centi\meter}   & dispersión de dos lecturas \\
Masa del trozo, $m$         & \SI{0,060 \pm 0,003}{\gram}       & resolución de la báscula \\
Longitud del trozo, $\ell$  & \SI{49,00 \pm 0,03}{\centi\meter} & resolución de la regla \\
\bottomrule
\end{tabular}
\end{table}

\subsection{Tensiones de los modos}\label{sec:tensiones}
La \cref{tab:datos} recoge las tensiones a las que se formó cada modo en las
cuatro configuraciones, junto con $\lambda_n^{2}=(2L/n)^{2}$. Con la cuerda
doble en $L_1$ solo pudimos distinguir con claridad los cuatro primeros
modos; con la doble en $L_2$, en cambio, llegamos hasta seis vientres. Como
predice la \cref{fig:modos}, la tensión cae rápidamente con el número de
vientres: del modo $n=1$ al $n=2$ se reduce aproximadamente a la cuarta
parte, y la cuerda doble necesita, para el mismo modo, aproximadamente el
doble de tensión que la simple.

\begin{table}[ht]
\centering
\caption{Tensión $T$ a la que aparece cada modo y cuadrado de la longitud de
onda, $\lambda_n^{2}=(2L/n)^{2}$, con su incertidumbre típica.}
\label{tab:datos}
\renewcommand{\arraystretch}{1.2}
\begin{tabular}{@{}c c c c @{\hspace{1.2em}} c c c@{}}
\toprule
 & \multicolumn{3}{c}{\scshape\color{azulprof}$L_2=\SI{38,8}{\centi\meter}$} &
   \multicolumn{3}{c@{}}{\scshape\color{azulprof}$L_1=\SI{51,9}{\centi\meter}$} \\
\cmidrule(lr){2-4}\cmidrule(l){5-7}
$n$ & $\lambda_n^{2}$ (\si{\meter\squared}) & $T$ simple (\si{\newton}) & $T$ doble (\si{\newton}) &
      $\lambda_n^{2}$ (\si{\meter\squared}) & $T$ simple (\si{\newton}) & $T$ doble (\si{\newton}) \\
\midrule
1 & \num{0,602 \pm 0,019}   & \num{0,15}  & \num{0,29}  & \num{1,08 \pm 0,04}    & \num{0,25}  & \num{0,5}   \\
2 & \num{0,151 \pm 0,005}   & \num{0,04}  & \num{0,08}  & \num{0,269 \pm 0,011}  & \num{0,07}  & \num{0,14}  \\
3 & \num{0,0669 \pm 0,0021} & \num{0,02}  & \num{0,04}  & \num{0,120 \pm 0,005}  & \num{0,03}  & \num{0,06}  \\
4 & \num{0,0376 \pm 0,0012} & \num{0,01}  & \num{0,02}  & \num{0,067 \pm 0,003}  & \num{0,008} & \num{0,017} \\
5 & \num{0,0241 \pm 0,0007} & \num{0,006} & \num{0,014} & \num{0,0431 \pm 0,0017}& \num{0,002} & ---         \\
6 & \num{0,0167 \pm 0,0005} & ---         & \num{0,01}  & ---                    & ---         & ---         \\
\bottomrule
\end{tabular}
\end{table}

\subsection{Ajustes lineales}\label{sec:ajustes}
La \cref{fig:ajustes} representa $\lambda^{2}$ frente a $T$ para las cuatro
series, y la \cref{tab:ajustes} recoge los parámetros de los ajustes por
mínimos cuadrados (\cref{met:regresion}). Las barras de error son las
incertidumbres típicas de $\lambda^{2}$ de la \cref{tab:datos}; como los
ajustes no están ponderados, no influyen en los parámetros.

\begin{figure}[ht]
\centering
\begin{tikzpicture}
  \pgfplotstableread{
    T      y       u
    0.15   0.6022  0.019
    0.04   0.1505  0.005
    0.02   0.0669  0.0021
    0.01   0.0376  0.0012
    0.006  0.0241  0.0007
  }\simpleLdos
  \pgfplotstableread{
    T      y       u
    0.25   1.0774  0.044
    0.07   0.2694  0.011
    0.03   0.1197  0.005
    0.008  0.0673  0.003
    0.002  0.0431  0.0017
  }\simpleLuno
  \pgfplotstableread{
    T      y       u
    0.29   0.6022  0.019
    0.08   0.1505  0.005
    0.04   0.0669  0.0021
    0.02   0.0376  0.0012
    0.014  0.0241  0.0007
    0.01   0.0167  0.0005
  }\dobleLdos
  \pgfplotstableread{
    T      y       u
    0.5    1.0774  0.044
    0.14   0.2694  0.011
    0.06   0.1197  0.005
    0.017  0.0673  0.003
  }\dobleLuno
  \begin{groupplot}[group style={group size=2 by 2, horizontal sep=1.1cm,
                                 vertical sep=0.35cm, x descriptions at=edge bottom},
                    datos, width=0.505\textwidth,
                    xticklabel style={/pgf/number format/fixed, /pgf/number format/precision=2},
                    yticklabel style={/pgf/number format/fixed, /pgf/number format/precision=2}]
    % ---- cuerda simple ----
    \nextgroupplot[height=6.2cm, xmin=0, xmax=0.27, ymin=0, ymax=1.2,
                   ylabel={$\lambda^{2}$ (\si{\meter\squared})},
                   title={\small\scshape\color{azulprof}cuerda simple},
                   legend style={at={(0.03,0.97)}, anchor=north west}]
      \addplot[ajuste, domain=0:0.27, samples=2] {4.0476*x-0.00668};
      \addplot[ajusteB, domain=0:0.27, samples=2] {4.2144*x+0.01196};
      \addplot[puntoA] table[x=T, y=y, y error=u] {\simpleLdos};
      \addplot[puntoB] table[x=T, y=y, y error=u] {\simpleLuno};
      \legend{ajuste $L_2$, ajuste $L_1$, datos $L_2$, datos $L_1$}
    % ---- cuerda doble ----
    \nextgroupplot[height=6.2cm, xmin=0, xmax=0.54, ymin=0, ymax=1.2,
                   title={\small\scshape\color{azulprof}cuerda doble},
                   legend style={at={(0.03,0.97)}, anchor=north west}]
      \addplot[ajuste, domain=0:0.54, samples=2] {2.0987*x-0.00912};
      \addplot[ajusteB, domain=0:0.54, samples=2] {2.1371*x+0.00038};
      \addplot[puntoA] table[x=T, y=y, y error=u] {\dobleLdos};
      \addplot[puntoB] table[x=T, y=y, y error=u] {\dobleLuno};
      \legend{ajuste $L_2$, ajuste $L_1$, datos $L_2$, datos $L_1$}
    % ---- residuos ----
    \nextgroupplot[height=3.2cm, xmin=0, xmax=0.27, ymin=-0.05, ymax=0.05,
                   ytick={-0.04,0,0.04},
                   ylabel={residuo (\si{\meter\squared})}, xlabel={$T$ (\si{\newton})}]
      \addplot[grispizarra, densely dashed, line width=0.6pt, domain=0:0.27, samples=2] {0};
      \addplot[puntoA] table[x=T, y expr=\thisrow{y}-(4.0476*\thisrow{T}-0.00668), y error=u] {\simpleLdos};
      \addplot[puntoB] table[x=T, y expr=\thisrow{y}-(4.2144*\thisrow{T}+0.01196), y error=u] {\simpleLuno};
    \nextgroupplot[height=3.2cm, xmin=0, xmax=0.54, ymin=-0.05, ymax=0.05,
                   ytick={-0.04,0,0.04}, xlabel={$T$ (\si{\newton})}]
      \addplot[grispizarra, densely dashed, line width=0.6pt, domain=0:0.54, samples=2] {0};
      \addplot[puntoA] table[x=T, y expr=\thisrow{y}-(2.0987*\thisrow{T}-0.00912), y error=u] {\dobleLdos};
      \addplot[puntoB] table[x=T, y expr=\thisrow{y}-(2.1371*\thisrow{T}+0.00038), y error=u] {\dobleLuno};
  \end{groupplot}
\end{tikzpicture}
\caption{$\lambda^{2}$ frente a $T$ para la cuerda simple (izquierda) y la
doble (derecha), en las dos distancias entre apoyos. Círculos: $L_2$;
cuadrados: $L_1$. Arriba, datos y rectas de regresión por mínimos cuadrados
(continua para $L_2$, discontinua para $L_1$); abajo, residuos. Las barras son
la incertidumbre típica $u(\lambda^{2})$. Obsérvese la diferencia de escala
horizontal: la cuerda doble requiere el doble de tensión.}
\label{fig:ajustes}
\end{figure}

\begin{table}[ht]
\centering
\caption{Parámetros de los ajustes $\lambda^{2}=a\,T+b$: valor con su
incertidumbre típica y, entre corchetes, la ampliada al \SI{95}{\percent}
($\nu=N-2$).}
\label{tab:ajustes}
\renewcommand{\arraystretch}{1.25}
\begin{tabular}{@{}l c c c c@{}}
\toprule
Serie & $N$ & $a$ (\si{\meter\squared\per\newton}) & $b$ (\si{\meter\squared}) & $r$ \\
\midrule
Simple, $L_2$ & 5 & \num{4,05 \pm 0,06}\ [\num{0,18}] & \num{-0,007 \pm 0,004}\ [\num{0,013}] & \num{0,9997} \\
Simple, $L_1$ & 5 & \num{4,21 \pm 0,15}\ [\num{0,5}]  & \num{0,012 \pm 0,018}\ [\num{0,06}]   & \num{0,998}  \\
Doble, $L_2$  & 6 & \num{2,10 \pm 0,03}\ [\num{0,08}] & \num{-0,009 \pm 0,004}\ [\num{0,010}] & \num{0,9996} \\
Doble, $L_1$  & 4 & \num{2,14 \pm 0,08}\ [\num{0,4}]  & \num{0,00 \pm 0,02}\ [\num{0,09}]     & \num{0,998}  \\
\bottomrule
\end{tabular}
\end{table}

Las cuatro series son claramente lineales ($r\geq\num{0,998}$) y las cuatro
ordenadas son compatibles con cero al \SI{95}{\percent}, como exige la
\cref{eq:lineal}. Los residuos, sin embargo, no son del todo aleatorios: en
las cuatro series los modos $n=2$ y $n=3$ quedan por debajo de la recta y los
de menor tensión, por encima. Es lo que cabría esperar si las tensiones más
pequeñas se leyeran algo por defecto, justo en la zona de la escala del
dinamómetro donde la lectura es menos fiable (\cref{obs:lecturas}).

\needspace{8\baselineskip}
\subsection{Comparación de las pendientes}\label{sec:comparacion}
\begin{estimacion}[Misma distancia, distinta densidad]\label{est:densidad}
Según la \cref{obs:pendientes}, el cociente entre la pendiente de la cuerda
simple y la de la doble debe valer 2. Propagando las incertidumbres de las
pendientes,
\begin{equation}\label{eq:cocientes}
  \frac{a_{\text{simple}}}{a_{\text{doble}}}\bigg|_{L_2}=\num{1,93 \pm 0,04},
  \qquad
  \frac{a_{\text{simple}}}{a_{\text{doble}}}\bigg|_{L_1}=\num{1,97 \pm 0,10}.
\end{equation}
En $L_1$ el cociente es plenamente compatible con 2; en $L_2$ se queda a
algo menos de dos incertidumbres típicas, en el límite de la compatibilidad
al \SI{95}{\percent}. Doblar la cuerda reduce la pendiente a la mitad, como
predice la teoría.
\end{estimacion}

\begin{estimacion}[Misma densidad, distinta distancia]\label{est:distancia}
Con la cuerda simple, las pendientes en $L_2$ y $L_1$ (\num{4,05} y
\num{4,21}) difieren en \num{0,17}, frente a una incertidumbre típica
combinada de \num{0,16}; con la doble (\num{2,10} y \num{2,14}), en
\num{0,04}, frente a \num{0,09}. En ambos casos las pendientes son
compatibles: cambiar la distancia entre apoyos no altera la pendiente, y
los puntos de ambas distancias caen sobre la misma recta.
\end{estimacion}

\subsection{Densidad lineal de la cuerda}\label{sec:densidad}
De cada pendiente obtenemos la densidad lineal con la \cref{eq:lineal}, y su
incertidumbre con la \cref{eq:umu}. Para compararlas entre sí y con la de la
báscula, dividimos por dos las de la cuerda doble (\cref{tab:densidad,fig:densidad}).

\begin{table}[ht]
\centering
\caption{Densidad lineal deducida de cada pendiente: valor con su
incertidumbre típica y, entre corchetes, la ampliada al \SI{95}{\percent}. En
la última columna, la densidad por hebra.}
\label{tab:densidad}
\renewcommand{\arraystretch}{1.25}
\begin{tabular}{@{}l c c@{}}
\toprule
Serie & $\mus=1/(af^{2})$ (\si{\kilogram\per\meter}) & Por hebra (\si{\kilogram\per\meter}) \\
\midrule
Simple, $L_2$ & $\num{1,34 \pm 0,03 e-4}$\ [\num{0,06e-4}] & $\num{1,34e-4}$ \\
Simple, $L_1$ & $\num{1,28 \pm 0,05 e-4}$\ [\num{0,14e-4}] & $\num{1,28e-4}$ \\
Doble, $L_2$  & $\num{2,58 \pm 0,05 e-4}$\ [\num{0,11e-4}] & $\num{1,29e-4}$ \\
Doble, $L_1$  & $\num{2,53 \pm 0,10 e-4}$\ [\num{0,4e-4}]  & $\num{1,27e-4}$ \\
\midrule
Media ponderada (por hebra) & & $\num{1,31 \pm 0,02 e-4}$ \\
Báscula, $\mub$             & & $\num{1,22 \pm 0,06 e-4}$ \\
\bottomrule
\end{tabular}
\end{table}

\begin{figure}[ht]
\centering
\begin{tikzpicture}
  \begin{axis}[datos, width=0.82\textwidth, height=5.4cm,
               xmin=1.0, xmax=1.55, ymin=0.4, ymax=6.6,
               ytick={1,...,6},
               yticklabels={báscula, media ponderada, doble $L_1$, doble $L_2$, simple $L_1$, simple $L_2$},
               yticklabel style={font=\small}, grid=major,
               xlabel={densidad lineal por hebra ($10^{-4}$ \si{\kilogram\per\meter})},
               xticklabel style={/pgf/number format/fixed, /pgf/number format/precision=2}]
    % banda de la báscula (95 %)
    \fill[acento!12] (axis cs:1.109,0.4) rectangle (axis cs:1.339,6.6);
    \draw[acento, line width=0.8pt, densely dashed] (axis cs:1.224,0.4) -- (axis cs:1.224,6.6);
    \addplot[only marks, mark=*, mark size=2.2pt, azulprof,
             error bars/.cd, x dir=both, x explicit,
             error bar style={line width=0.8pt, azulprof},
             error mark options={rotate=90, mark size=3pt, line width=0.8pt, azulprof}]
      coordinates {(1.336,6) +- (0.057,0) (1.283,5) +- (0.137,0)
                   (1.289,4) +- (0.053,0) (1.265,3) +- (0.186,0)};
    \addplot[only marks, mark=diamond*, mark size=3.2pt, azulacero,
             error bars/.cd, x dir=both, x explicit,
             error bar style={line width=0.8pt, azulacero},
             error mark options={rotate=90, mark size=3pt, line width=0.8pt, azulacero}]
      coordinates {(1.307,2) +- (0.043,0)};
    \addplot[only marks, mark=square*, mark size=2.2pt, acento,
             error bars/.cd, x dir=both, x explicit,
             error bar style={line width=0.8pt, acento},
             error mark options={rotate=90, mark size=3pt, line width=0.8pt, acento}]
      coordinates {(1.224,1) +- (0.115,0)};
  \end{axis}
\end{tikzpicture}
\caption{Densidad lineal por hebra obtenida de cada ajuste, su media
ponderada y el valor de la báscula, con sus incertidumbres ampliadas al
\SI{95}{\percent}. La banda sombreada es el intervalo al \SI{95}{\percent}
de la báscula.}
\label{fig:densidad}
\end{figure}

Las cuatro estimaciones son compatibles entre sí. Combinándolas con una media
ponderada por el inverso del cuadrado de las incertidumbres de las
pendientes, y añadiendo después la contribución de la frecuencia, común a
todas, obtenemos
\begin{destacado}[Densidad lineal a partir de las pendientes]
\vspace{-6pt}
\begin{equation}\label{eq:mumedia}
  \mus=\num{1,31 \pm 0,02 e-4}\ \si{\kilogram\per\meter}
  \quad(\text{incertidumbre típica}),
\end{equation}
{\small frente a $\mub=\num{1,22 \pm 0,06 e-4}\ \si{\kilogram\per\meter}$
de la báscula: una diferencia del \SI{7}{\percent}, de unas $\num{1,3}$
incertidumbres típicas combinadas, por lo que ambos valores son compatibles.}
\end{destacado}

% =========================================================================
\newpage
\seccion{Discusión y conclusiones}{Lo que aprendemos}\label{sec:discusion}
% =========================================================================
\begin{conclusion}[Linealidad y modos]\label{conc:lineal}
En las cuatro configuraciones, $\lambda^{2}$ crece linealmente con la tensión
($r\geq\num{0,998}$) y las ordenadas son compatibles con cero: las
longitudes de onda que impone la geometría, $\lambda_n=2L/n$, y la velocidad
de propagación $v=\sqrt{T/\mu}$ describen bien la cuerda. La tensión de cada
modo cae como $1/n^{2}$, tal y como se espera a frecuencia fija.
\end{conclusion}

\begin{conclusion}[Comparación de las pendientes]\label{conc:pendientes}
Para una misma cuerda, la pendiente no depende de la distancia entre apoyos
(\num{4,05} y \num{4,21} con la simple; \num{2,10} y \num{2,14} con la doble,
en \si{\meter\squared\per\newton}). Para una misma distancia, doblar la cuerda
reduce la pendiente a la mitad: los cocientes valen \num{1,93 \pm 0,04} y
\num{1,97 \pm 0,10}, compatibles con 2.
\end{conclusion}

\begin{conclusion}[Densidad lineal]\label{conc:densidad}
Las pendientes dan una densidad por hebra de
$\num{1,31 \pm 0,02 e-4}\ \si{\kilogram\per\meter}$, compatible con la de la
báscula, $\num{1,22 \pm 0,06 e-4}\ \si{\kilogram\per\meter}$, aunque algo
mayor. El método dinámico resulta, de hecho, más preciso que la pesada, cuya
incertidumbre está dominada por la resolución de \SI{0,01}{\gram} sobre una
masa de solo \SI{0,06}{\gram}.
\end{conclusion}

\begin{observacion}[Fuentes de error]\label{obs:errores}
La diferencia sistemática entre ambas densidades y el patrón de los residuos
apuntan a tres limitaciones: la lectura de las tensiones más pequeñas, en el
límite de la escala del dinamómetro; la indefinición de la longitud que
vibra, cuyas dos lecturas difieren hasta en \SI{2}{\centi\meter}, y la
pesada de una masa tan pequeña. Pesar un trozo de cuerda más largo y emplear
un dinamómetro de menor rango en los modos altos reduciría las tres.
\end{observacion}

En conjunto, la práctica cumple sus objetivos: con una regla, un dinamómetro
y un estroboscopio hemos comprobado la cuantización de los modos de una
cuerda y medido su densidad lineal con más precisión que pesándola.

% =========================================================================
\newpage
%                              BIBLIOGRAFÍA
% =========================================================================
\printbibliography[title={Referencias}]

\end{document}
